\section*{Chapter \ref{ch:qm:multivariate-series-and-cointegration} --- Multivariate Series and Cointegration}

\begin{solution}{exo:qm:multivariate-series-and-cointegration:1}
The eigenvalues of the first matrix are $0.9$ and $0.1$: stationary. Those of the second are $1.0$ and $0.2$: a \omterm{def:qm:linear-time-series:unitroot}{unit root}, so not stationary; the sum $x + y$
follows a random walk while the difference $x - y$ is stationary (with coefficient 0.2): the two series are cointegrated with $\beta = (1, -1)$.
\end{solution}

\begin{solution}{exo:qm:multivariate-series-and-cointegration:2}
$k = 2$, $r = 1$ (the vector $(1, -0.9)$ up to scale), and $k - r = 1$ common trend.
\end{solution}

\begin{solution}{exo:qm:multivariate-series-and-cointegration:3}
$-3.74066 - 8.5631/500 - 10.852/500^2 + 27.982/500^3 = -3.758$.
\end{solution}

\begin{solution}{exo:qm:multivariate-series-and-cointegration:4}
$\Delta y_t = -0.5(y_{t-1} - 0.9x_{t-1}) + e_t$ and $\Delta x_t = u_t$. So $\beta = (1, -0.9)^\top$ and $\alpha = (-0.5, 0)^\top$ for $(y, x)$: $y$ corrects half of any deviation
each period; $x$ does not correct at all (it is weakly exogenous).
\end{solution}

\begin{solution}{exo:qm:multivariate-series-and-cointegration:5}
The Cholesky factor writes the second residual as $0.9z_1 + \sqrt{0.19}\,z_2$, so 81\% of its one-step variance is attributed to the first shock; with the opposite ordering
it would be 19\% (or 0\% of the first variable's to the second shock in the original ordering).
\end{solution}

\begin{solution}{exo:qm:multivariate-series-and-cointegration:6}
Trace for $r = 0$: $-1\,000(\ln0.95 + \ln0.99 + \ln0.998) = 63.35$; $r \le 1$: 12.05; $r \le 2$: 2.00. Against 35.23, 20.36 and 9.23: reject $r = 0$, not $r \le 1$, so rank one.
\end{solution}

\begin{solution}{exo:qm:multivariate-series-and-cointegration:7}
\texttt{rolling\_weights()}: no \omterm{def:qm:multivariate-series-and-cointegration:coint}{cointegration} at 5\% in 20 of the 46 windows; where one vector is found, the 2-year weight ranges from 0.27 to 0.63 and the 10-year weight
from 0.40 to 0.97.
\end{solution}

\begin{solution}{exo:qm:multivariate-series-and-cointegration:8}
The residual of an estimated regression must be judged with Engle--Granger critical values, $-3.74$ at 5\% for three variables, not $-2.86$; $-3.1$ does not reject
the null of no \omterm{def:qm:multivariate-series-and-cointegration:coint}{cointegration}. With the Dickey--Fuller value, independent random walks would pass 27\% of the time. The weights should also be checked for stability before
the fly is traded.
\end{solution}

\begin{solution}{pb:qm:multivariate-series-and-cointegration:1}
\textbf{1.} Dickey--Fuller statistics $-1.30$, $-1.33$, $-1.31$: no yield rejects a \omterm{def:qm:linear-time-series:unitroot}{unit root}.
\textbf{2.} AIC keeps improving up to the maximum of ten lags; the VAR(10) of the changes is stable.
\textbf{3.} 7.9, 7.6 and 7.1 bp; correlations 0.90 (2y--5y), 0.82 (2y--10y), 0.94 (5y--10y).
\textbf{4.} With the 2-year ordered first, a 7.9 bp shock moves the 2-, 5- and 10-year yields by 10.1, 8.3 and 6.7 bp cumulatively over ten days, and explains 67\% of the
10-year's ten-day forecast-error variance.
\textbf{5.} All six directions are significant at 5\% ($F$ from 1.88 to 7.16 with 10 and 12\,532 degrees of freedom): with so much data, small lead-lag effects from
timing and common news are detectable; they are predictive precedence, not causation.
\textbf{6.} 59.77, 14.93 and 2.34 against 35.23, 20.36 and 9.23: rank one.
\textbf{7.} $0.41 \times$ 2-year $- 1 \times$ 5-year $+ 0.61 \times$ 10-year: the wings sum to 1.02 (level-neutral), and the tilt toward the 10-year also neutralises most of the
slope.
\textbf{8.} The 5-year on the others: $0.40$ and $0.62$, residual statistic $-10.04$ against a 5\% critical value of $-3.74$.
\textbf{9.} 43 trading days (AR coefficient 0.9841) against 63 (0.9891).
\textbf{10.} 2.06 and 2.15 bp a day (the one-two-one fly scaled to the same 5-year weight).
\textbf{11.} In 26 of 46 windows (20 find none at 5\%).
\textbf{12.} The 2-year weight from 0.27 to 0.63 and the 10-year from 0.40 to 0.97, among the windows that find one vector.
\textbf{13.} Least squares chooses the combination that looks most stationary in the sample, so the residual's statistic is pushed left even when nothing is
cointegrated; the more variables, the further.
\textbf{14.} Apply Kendall's correction $(1 + 3\hat\rho)/n$ to the fly's AR(1) coefficient: 0.9841 becomes 0.9844 with 12\,574 days and the \omterm{def:qm:stochastic-differential-equations:ou}{half-life} 44 days
instead of 43; then simulate its distribution as in chapter~17.
\textbf{15.} That the weights of the next years differ, so the traded fly carries level or slope exposure it was meant to remove, and reverts more slowly than the
backtest promised.
\textbf{16.} The Johansen fly, with weights re-estimated on a recent window and bounded near the long-run values, cross-checked against duration-neutral weights.
\textbf{17.} By the fly's stationary standard deviation, $2.06/\sqrt{1 - 0.9841^2} = 11.6$ bp: size so that a two- or three-standard-deviation
move (25 to 35 bp) is survivable, and expect a typical deviation to halve in about two months.
\textbf{18.} Because the curve's factor loadings change with the rate regime; re-estimate on a rolling window every month or quarter and trade only when the rank test
still finds the vector.
\textbf{19.} Named result: \emph{the fly that mean-reverts}: the Johansen cointegrating fly of the 2-, 5- and 10-year yields, 1976--2026, weights the wings $0.41$ and $0.61$ per
unit of the 5-year and reverts with a \omterm{def:qm:stochastic-differential-equations:ou}{half-life} of 43 trading days, against 63 for the one-two-one fly.
\textbf{20.} It promises that a basket's deviations are temporary on average; it does not promise how long they last, how large they get first, or that the weights will
stay the same.
\end{solution}

\begin{solution}{iq:qm:multivariate-series-and-cointegration:1}
Correlation describes co-movement of changes over a horizon; \omterm{def:qm:multivariate-series-and-cointegration:coint}{cointegration} says that the levels share stochastic trends, so some combination of them is stationary. Two
series can be highly correlated and drift apart forever, or weakly correlated day to day and tied together in the long run.
\iqlookfor{Changes versus levels, and a counterexample either way.}
\end{solution}

\begin{solution}{iq:qm:multivariate-series-and-cointegration:2}
Choose weights that make the fly stationary and fast-reverting: estimate the \omterm{def:qm:multivariate-series-and-cointegration:coint}{cointegrating vector} (Johansen or Engle--Granger with the right critical values),
compare with duration- or PCA-neutral weights, check stability on rolling windows, and measure the \omterm{def:qm:stochastic-differential-equations:ou}{half-life} and volatility of the resulting fly.
\iqlookfor{Stationarity as the criterion, the right test, and stability.}
\end{solution}

\begin{solution}{iq:qm:multivariate-series-and-cointegration:3}
The coefficients were estimated to make the residual as small, and so as stationary-looking, as possible; under the null the residual's statistic is shifted left, and the
shift grows with the number of regressors. Use Engle--Granger (MacKinnon) critical values: $-3.34$ for two variables, $-3.74$ for three at 5\%, asymptotically.
\iqlookfor{The pre-fitting argument and the right table.}
\end{solution}

\begin{solution}{iq:qm:multivariate-series-and-cointegration:4}
Whether the past of one series improves forecasts of another given its own past; it tests predictive precedence within the information set used. It does not establish
causation: omitted common drivers, timing differences in data, and expectations can all produce it.
\iqlookfor{Prediction, not causation, and the omitted-variable caveat.}
\end{solution}

\begin{solution}{iq:qm:multivariate-series-and-cointegration:5}
Because correlated shocks are split by a Cholesky factorisation: the first variable gets all the common part of the shock. Different orderings attribute the
contemporaneous correlation differently; report the ordering, try alternatives, or use generalised responses.
\iqlookfor{Orthogonalisation and its arbitrariness.}
\end{solution}

\begin{solution}{iq:qm:multivariate-series-and-cointegration:6}
After partialling out the short-run lags, compute the canonical correlations between the differences $\Delta Y_t$ and the lagged levels $Y_{t-1}$. Stationary
combinations of the levels predict the differences, so they have large canonical correlations; the number of significantly nonzero ones is the rank, the
eigenvectors estimate $\beta$, and the trace statistic sums $-T\ln(1 - \lambda_i)$ over the smallest ones.
\iqlookfor{Reduced-rank regression through canonical correlations.}
\end{solution}
